\section{Literature position}\label{sec:literature}

The closest located literature falls into four groups.

\emph{Persistent homology of random fields.} Adler--Bobrowski--Borman--Subag--Weinberger \cite{ABBSW} and Pranav \cite{Pranav} study persistent homology of random and Gaussian fields, principally through general structure and computation; Bobrowski--Borman \cite{BobrowskiBorman} compute expected barcode descriptors of Gaussian fields (alternating sums of bar lengths over homological degrees), which integrate lifetimes rather than resolving the $H_0$ lifetime measure near zero. Feldbrugge, van Engelen, van de Weygaert, Pranav and Vegter \cite{Feldbrugge} express expected two-dimensional superlevel persistence diagrams through candidate maximum--saddle densities and the probability that the saddle kills the component born at the given maximum: this is the selection-kernel architecture of \S\ref{sec:proof}, without a short-gap asymptotic. Chazal--Divol \cite{ChazalDivol} prove absolute continuity of expected persistence diagrams for random point-cloud filtrations and treat Brownian sublevel sets; their results justify the existence of such densities in those settings and do not assert a general near-diagonal rate for smooth stationary Gaussian fields.

\emph{Critical points of Gaussian fields.} Klein--Agam \cite{KleinAgam} describe the short-distance correlation of critical points, of order $r^{2-d}$, and Azaïs--Delmas \cite{AzaisDelmas} prove it, with explicit constants, for critical points of adjacent indices; in polar coordinates this is the $r\,dr$ pair measure behind \eqref{eq:radial}. Ladgham--Lachièze-Rey \cite{LadghamLR} prove local repulsion of planar critical points, and Ancona--Gass--Letendre--Stecconi \cite{AGLS} analyze Kac--Rice singularities and cumulants. None of these imposes a height gap or the elder pairing, and the cubic fold exponent itself is classical.

\emph{Limit theorems for topological functionals and short features.} Hirsch--Lachièze-Rey \cite{HirschLR} prove functional central limit theorems for topological functionals of Gaussian fields on growing domains. Chenavier--Hirsch \cite{ChenavierHirsch} prove Poisson limits for persistent cycles of random point clouds with extremal (minimal and maximal) lifetimes, and Hirsch--Lundbye--Otto \cite{HLO} for large-lifetime cycles, including on the torus. Perez \cite{Perez} obtains small-cutoff asymptotics of bar counts for continuous semimartingales, governed by quadratic variation; that concerns rough paths and counts above a cutoff, a different regularity class.

\emph{Universality of persistence ratios.} For random geometric complexes, Bobrowski and Skraba conjectured that suitably centered distributions of $\log\log$ of the death/birth ratio are universal \cite{BobrowskiSkraba2023}, and proved a weak form at fixed dimension and positive homological degree \cite{BobrowskiSkraba2024}. Heo \cite{Heo2026} computes a second-order near-diagonal coefficient of the ratio law for degree-one \v{C}ech persistence in dimensions two and three, finds different values, and concludes that the strong form fails. These results concern point-cloud filtrations, positive homological degree and a multiplicative statistic; the present law concerns the additive lifetime of $H_0$ bars of a smooth Gaussian field.

\emph{Dimension one.} Rice \cite{Rice} and Cartwright--Longuet-Higgins \cite{CLH} studied the heights of the maxima of a stationary Gaussian process, the latter also crest-to-trough heights for narrow spectra. The amplitude between a local maximum and the next minimum was studied by Lindgren \cite{Lindgren1972}, using Slepian-model conditional formulas, and by Lindgren--Rychlik \cite{LindgrenRychlik}; Lindgren \cite{Lindgren1970} shows that near a maximum of very low height the process is a random quartic, a low-level regime rather than a small-amplitude law. The exact joint law of adjacent extremes \cite{Lindgren2019} contains the small-amplitude behaviour of \S\ref{ssec:d1} implicitly, but we have not found the $-1/3$ law stated there.

We have not located, in these sources or in the indexed and web searches recorded in \rec{lit} (the latest on 13 October 2026), a prior theorem giving the fixed-side three-dimensional near-diagonal $H_0$ density \eqref{eq:main}, the coefficient reduction \eqref{eq:c} or the correction \eqref{eq:correction}, nor any near-diagonal lifetime law for the $H_0$ persistence of smooth Gaussian fields in dimension at least two. Substantial constituent mechanisms --- the fold, the close-pair correlation, the selection kernel --- are prior art. This is evidence of novelty, not a substitute for a systematic MathSciNet and zbMATH search with backward and forward citation chasing, which remains to be done before submission.

\emph{Relation to the parent record.} The leading law \eqref{eq:main} and the formula \eqref{eq:c} are the case $d=3$, $L=24$ of Theorem~C and (15.2) of the parent record \rec{parent}, which holds for every $d\ge2$ and $L>0$; the selection bound \eqref{eq:selection} is its Theorem~A, and Proposition~\ref{prop:far} is its (14.1). Specific to the side-$24$ field are the closed form \eqref{eq:cinf}, the correction \eqref{eq:correction} with its error bound, the certified constants of Theorem~\ref{thm:three-term}(iii), and the second, three-dimensional proof of \eqref{eq:selection} in Supplement~\S S5. The parent record is not a publication. The Supplement proves the results behind \S\ref{sec:proof} and Theorem~\ref{thm:three-term}, but imports some statements from the parent record, listed in its Sections S7, S8 and S10, among them the cap theorem, the marked Kac--Rice identity, the far bound (14.1) and the coefficient definition.
